% ===== PRÉAMBULE CANONIQUE — à coller VERBATIM en tête de CHAQUE .tex =====
\documentclass[11pt,a4paper]{article}
\usepackage[utf8]{inputenc}\usepackage[T1]{fontenc}\usepackage{lmodern}
\usepackage[french]{babel}
\usepackage{amsmath,amssymb,mathtools}\usepackage{icomma}
\usepackage{siunitx}\sisetup{locale=FR}  % \qty{}{}, \si{}, \ang{}
\usepackage{xcolor}\usepackage{colortbl}
\usepackage{tikz}\usepackage{pgfplots}\pgfplotsset{compat=1.18}
\usepackage[siunitx,europeanresistors]{circuitikz} % circuits (élec.)
\usepackage[most]{tcolorbox}\usepackage{enumitem}\usepackage{array,booktabs}
\usepackage[a4paper,margin=2.2cm]{geometry}
\usetikzlibrary{arrows.meta,calc,angles,quotes,positioning,patterns,%
  backgrounds,shapes.geometric,decorations.pathmorphing,decorations.markings,intersections}
\definecolor{primary}{RGB}{222,49,99}\definecolor{primarydark}{RGB}{150,22,55}
\definecolor{defcol}{RGB}{0,114,178}\definecolor{methcol}{RGB}{52,92,148}
\definecolor{excol}{RGB}{0,135,90}\definecolor{attncol}{RGB}{200,40,55}
\tcbset{boxbase/.style={enhanced,breakable,boxrule=0.9pt,arc=2.5pt,
  left=3mm,right=3mm,top=2.3mm,bottom=2.3mm,fonttitle=\bfseries,coltitle=white}}
% --- boîtes TUTORIEL (noms EXACTS, ne pas renommer) ---
\newtcolorbox{cle}[1][]{boxbase,colback=defcol!6,colframe=defcol,
  title={\textsc{À retenir}\ifx\relax#1\relax\relax\else\textmd{ : \itshape #1}\fi}}
\newtcolorbox{methode}[1][]{boxbase,colback=methcol!7,colframe=methcol,sharp corners,
  title={\textsc{Méthode}\ifx\relax#1\relax\relax\else\textmd{ --- \itshape #1}\fi}}
\newtcolorbox{exemplebox}[1][]{boxbase,colback=excol!5,colframe=excol,
  title={\textsc{Exemple}\ifx\relax#1\relax\relax\else\textmd{ : \itshape #1}\fi}}
\newtcolorbox{piege}[1][]{boxbase,colback=attncol!6,colframe=attncol,
  title={\textsc{Piège}\ifx\relax#1\relax\relax\else\textmd{ : \itshape #1}\fi}}
\newtcolorbox{remarque}[1][]{boxbase,colback=black!4,colframe=black!45,
  title={\textsc{Remarque}\ifx\relax#1\relax\relax\else\textmd{ : \itshape #1}\fi}}
% --- boîtes EXERCICE / CORRIGÉ (noms EXACTS, auto-comptées) ---
\newtcolorbox[auto counter]{exo}[1][]{boxbase,colback=primary!4,colframe=primary,
  title={\textsc{Exercice}~\thetcbcounter\ifx\relax#1\relax\relax\else\textmd{ --- \itshape #1}\fi}}
\newtcolorbox[auto counter]{corrige}[1][]{boxbase,colback=excol!4,colframe=excol,
  title={\textsc{Corrigé}~\thetcbcounter\ifx\relax#1\relax\relax\else\textmd{ --- \itshape #1}\fi}}
\newcommand{\vv}[1]{\vec{#1}}\newcommand{\ds}{\displaystyle}
\newcommand{\R}{\mathbb{R}}\newcommand{\N}{\mathbb{N}}\newcommand{\Z}{\mathbb{Z}}
% (un \usepackage{fancyhdr}+en-tête est possible ; il est ignoré côté web)
% =========================================================================
\begin{document}
\sisetup{per-mode=symbol}

\begin{exo}[quel spectre est correct ?]
On propose trois représentations. Le vecteur $\vv{E}$ (en trait épais) est censé indiquer le champ.
\begin{center}
\begin{tikzpicture}[scale=0.72,>=Stealth,font=\scriptsize,
    ch/.style={circle,minimum size=6mm,inner sep=0pt,text=white,font=\bfseries\scriptsize}]
  % (A) E perpendiculaire a la ligne
  \begin{scope}[xshift=-5.2cm]
    \draw[excol,line width=0.9pt] (-1.3,-0.6) .. controls (-0.4,0.7) and (0.4,-0.7) .. (1.3,0.6);
    \draw[->,primary,line width=1.3pt] (0,0)--(0.6,0.6);
    \node[primary,font=\scriptsize] at (0.9,0.75){$\vv{E}$};
    \node[primarydark,font=\bfseries] at (0,-1.5){(A)};
  \end{scope}
  % (B) dipole inverse
  \begin{scope}[xshift=0cm]
    \node[ch,fill=attncol] (p) at (-1.1,0) {$+$};
    \node[ch,fill=defcol] (n) at (1.1,0) {$-$};
    \draw[->,primary,line width=1.1pt] (0.8,0)--(-0.8,0);
    \node[primarydark,font=\bfseries] at (0,-1.5){(B)};
  \end{scope}
  % (C) champ uniforme, E tangent
  \begin{scope}[xshift=5.2cm]
    \foreach \x in {-1,-0.4,0.2,0.8}{\draw[excol,line width=0.9pt] (\x,-0.9)--(\x,0.9);}
    \draw[->,primary,line width=1.3pt] (-0.4,-0.9)--(-0.4,0.9);
    \node[primary,font=\scriptsize] at (0.05,0.7){$\vv{E}$};
    \node[primarydark,font=\bfseries] at (-0.1,-1.5){(C)};
  \end{scope}
\end{tikzpicture}
\end{center}
\begin{enumerate}[label=\alph*)]
  \item Lequel de ces trois schémas représente correctement un champ électrique ?
  \item Explique l'erreur commise dans chacun des deux autres.
\end{enumerate}
\end{exo}

\begin{exo}[condensateur : attention au piège]
Un condensateur plan est formé de deux disques de rayon $R = \SI{10}{\centi\metre}$, séparés par
$d = 2R$, soumis à une tension $U_{AB} = \SI{100}{\volt}$ avec $V_A > V_B$.
\begin{enumerate}[label=\alph*)]
  \item Calcule l'intensité du champ uniforme entre les armatures. \emph{(Le rayon est-il utile ?)}
  \item En admettant que le potentiel varie linéairement, quelle est la différence de potentiel
        $U_{AP} = V_A - V_P$ entre l'armature $A$ et le point $P$ situé au \emph{milieu} des deux
        armatures ?
\end{enumerate}
\end{exo}

\begin{exo}[analyse complète d'un spectre]
Deux charges ponctuelles $A$ et $B$ sont reliées par des lignes de champ : \textbf{12} lignes
\emph{partent} de $A$, dont \textbf{6} \emph{aboutissent} à $B$ (les autres s'échappent vers
l'infini).
\begin{enumerate}[label=\alph*)]
  \item Quels sont les signes de $A$ et de $B$ ?
  \item Quelle charge est la plus grande, et dans quel rapport ?
  \item Que deviennent les $6$ lignes de $A$ qui n'aboutissent pas à $B$ ?
\end{enumerate}
\end{exo}

\begin{exo}[vrai ou faux : lignes de champ et condensateur]
Pour chaque proposition, réponds par \textbf{vrai} ou \textbf{faux}, puis donne le nombre total de
propositions correctes.
\begin{enumerate}[label=\alph*)]
  \item Les lignes de champ sont orientées dans le sens du vecteur champ électrique.
  \item Deux lignes de champ peuvent se croiser en un point où le champ est nul.
  \item Dans un condensateur, $\vv{E}$ est orienté vers l'armature de potentiel le plus \emph{haut}.
  \item Là où les lignes de champ sont resserrées, le champ est plus intense.
  \item Entre deux charges positives égales, il existe un point où le champ est nul.
\end{enumerate}
\end{exo}

\end{document}
